%% ma: L12-B16-0014
%% lop: 12
%% chuong: 5
%% bai: 16
%% dang: 12.16.11
%% loai: TN4
%% muc: VDC
%% dap_an: "B"
%% nguon: "Ảnh giáo viên gửi 04/11/2026 (Huỳnh Văn Ánh, Chuyên đề V - Phương pháp tọa độ trong không gian), Đề test theo bài, Phần I Câu 5"
%% kien_thuc: "Bài 16 (góc giữa hai đường thẳng), Bài 15 (giao điểm đường thẳng và mặt phẳng), Bài 14 (vectơ pháp tuyến), khoảng cách từ điểm đến đường thẳng (hình học phẳng, tam giác vuông), Lớp 11 Bài 23 (đường thẳng vuông góc với mặt phẳng)"
%% trang_thai: cho-duyet
%% ly_do: "Đề gốc in (P): x/1 + y/2 + z/2 = 1, tức 2x + y + z - 2 = 0. Với mặt phẳng này, (P) không vuông góc d và sin của góc giữa d1, d2 bằng 2 căn 6 chia 5 (khoảng 0,98), không có trong 4 phương án. Lời giải gốc dùng (P) vuông góc d và (P) cắt d tại I(1;0;0), tương ứng (P): x + 2y + 2z - 1 = 0 (vectơ pháp tuyến (1;2;2) trùng vectơ chỉ phương của d). Bản kho tạm sửa đề thành x + 2y + 2z - 1 = 0 cho khớp lời giải và đáp án B; cần giáo viên xác nhận đề đúng."
%% ghi_chu: "Với (P): x+2y+2z-1=0, A(1;-1;1) thuộc (P), sin phi = 1 (sympy), khớp đáp án gốc B. Hình gốc nằm trong lời giải và có đường phụ H, I (gợi ý lời giải) nên không đưa hình vào đề, theo quy ước hình không chứa gợi ý. Phân vân giữa VD và VDC, đề dựng hình tự do nên chọn VDC."
\begin{ex}
Trong không gian $Oxyz$, cho mặt phẳng $(P)\colon x+2y+2z-1=0$ và đường thẳng $d\colon\dfrac{x-2}{1}=\dfrac{y-2}{2}=\dfrac{z-2}{2}$. Biết rằng trong mặt phẳng $(P)$ có hai đường thẳng $d_1,d_2$ cùng đi qua điểm $A(1;-1;1)$ và cùng cách đường thẳng $d$ một khoảng bằng $1$. Tính $\sin\varphi$ với $\varphi$ là góc giữa hai đường thẳng $d_1,d_2$.
\choice{$\dfrac12$}{\True $1$}{$\dfrac{\sqrt3}{2}$}{$\dfrac37$}
\loigiai{$d$ có vectơ chỉ phương $\vec u=(1;2;2)$ trùng với vectơ pháp tuyến của $(P)$ nên $d\perp(P)$.

Thay $x=2+t$, $y=2+2t$, $z=2+2t$ vào $(P)$: $9t+9=0\Leftrightarrow t=-1$, nên $d\cap(P)=I(1;0;0)$ và $AI=\sqrt2$.

Vì $d_i\subset(P)$ và $d\perp(P)$ tại $I$ nên khoảng cách giữa $d_i$ và $d$ bằng khoảng cách từ $I$ đến $d_i$. Gọi $H$ là hình chiếu vuông góc của $I$ lên $d_1$, ta có $HI=1$.

Trong tam giác vuông $HAI$: $\sin\widehat{HAI}=\dfrac{HI}{AI}=\dfrac1{\sqrt2}\Rightarrow\widehat{HAI}=45^\circ$.

Hai đường thẳng $d_1,d_2$ đối xứng nhau qua $AI$ nên $\varphi=2\cdot45^\circ=90^\circ$ và $\sin\varphi=1$. Chọn B.}
\end{ex}
